\documentclass[../../main.tex]{subfiles}

\begin{document}

\section{关系}\label{section:ch01-关系}

从某种意义来说，关系是一个比函数更为广泛的概念。本节要给出数学家们所讲的关系的定义，并且研究在数学中经常出现的两种关系，即等价关系和全序关系。全序关系的概念在本书中自始至终都要用到，而等价关系的概念在第 22 节之前不会用到。

\begin{definition}[关系]\label{definition:ch01-关系}
集合 $A$ 上的一个关系（relation）是笛卡儿积 $A \times A$ 的一个子集 $C$。
\end{definition}

如果 $C$ 是 $A$ 中的一个关系，我们用记号 $xCy$ 表示 $(x, y) \in C$，读作“$x$ 与 $y$ 有关系 $C$”。

函数 $f: A \longrightarrow A$ 的一个指派法则 $r$ 也是 $A \times A$ 的一个子集，但它是相当特殊的一类子集，它使得 $A$ 的每一个元素作为 $r$ 的一个元素的第一个坐标恰好出现一次。$A \times A$ 的任意子集都是 $A$ 中的关系。

\begin{example}\label{example:ch01-人物关系}
设 $P$ 为全世界所有人的集合，$D \subset P \times P$ 定义为
\begin{align*}
D = \{(x, y) \mid x \text{ 是 } y \text{ 的一个后代}\},
\end{align*}
则 $D$ 是集合 $P$ 中的一个关系。“$x$ 与 $y$ 有关系 $D$”与“$x$ 是 $y$ 的后代”这两句话，说的是同一件事，即 $(x, y) \in D$。$P$ 中还有另外两个关系：
\begin{align*}
B = \{(x, y) \mid x \text{ 有一个祖先也是 } y \text{ 的祖先}\}.
\end{align*}
\begin{align*}
S = \{(x, y) \mid x \text{ 的父母就是 } y \text{ 的父母}\}.
\end{align*}
我们称 $B$ 为“血缘关系”，称 $S$ 为“同胞关系”。这三种关系有完全不同的性质。例如，血缘关系有对称性（如果 $x$ 与 $y$ 有血缘关系，则 $y$ 与 $x$ 有血缘关系），而后代关系则不然。我们马上要研究这些关系。
\end{example}

\subsection*{等价关系与分拆}

集合 $A$ 中的一个等价关系（equivalence relation）是 $A$ 上满足下面三条性质的一个关系 $C$：
\begin{enumerate}[(1)]
\item（自反性）对于 $A$ 中每一个 $x$，有 $xCx$.
\item（对称性）若 $xCy$，则 $yCx$.
\item（传递性）若 $xCy$ 和 $yCz$，则 $xCz$.
\end{enumerate}

\begin{example}\label{example:ch01-等价关系例}
在\cref{example:ch01-人物关系}所定义的关系中，后代关系 $D$ 既没有自反性，又没有对称性；而血缘关系 $B$ 没有传递性。（尽管我和我的妻子都与孩子有血缘关系，但是我们之间没有血缘关系！）然而可以验证同胞关系是一个等价关系。
\end{example}

虽然关系是一个集合，但并不一定要用大写字母或任何类型的字母来记关系。用另外的符号将更为合适。经常以表示一个等价关系的是“波纹”号 $\sim$。应用这个记号，等价关系的性质可写成：
\begin{enumerate}[(1)]
\item 对于 $A$ 中每一个 $x$，$x \sim x$.
\item 若 $x \sim y$，则 $y \sim x$.
\item 若 $x \sim y$ 和 $y \sim z$，则 $x \sim z$.
\end{enumerate}

有些特殊的等价关系使用了一些别的符号，它们在本书中许多地方都会遇到。

给定集合 $A$ 中的一个等价关系 $\sim$，和 $A$ 的一个元素 $x$，可以由下式定义 $A$ 的子集 $E$，称为由 $x$ 决定的等价类（equivalence class）：
\begin{align*}
E = \{y \mid y \sim x\}.
\end{align*}
注意，由 $x$ 决定的等价类 $E$ 包含 $x$，这是因为 $x \sim x$。等价类具有下列性质：

\begin{lemma}[等价类的性质]\label{lemma:ch01-等价类}
两个等价类 $E$ 和 $E'$ 或者无交或者相等。
\end{lemma}

\par
\begin{figure}[htbp]
\centering
\includegraphics[width=0.3\textwidth]{ch01-fig3-1}
\caption{两个等价类或者无交或者相等}\label{figure:ch01-等价类}
\end{figure}

\begin{proof}
设 $E$ 是由 $x$ 决定的等价类，$E'$ 是由 $x'$ 决定的等价类。若 $E \cap E'$ 不是空集，则有 $y \in E \cap E'$，见\cref{figure:ch01-等价类}。下面证明 $E = E'$。


根据定义可知 $y \sim x$ 和 $y \sim x'$。根据对称性可得 $x \sim y$ 和 $y \sim x'$。应用传递性则有 $x \sim x'$。若 $w$ 为 $E$ 的任意一点，根据定义有 $w \sim x$。再用一次传递性得到 $w \sim x'$。因此，$E \subset E'$。

可以完全对称地推出 $E' \subset E$。所以 $E = E'$。
\end{proof}

给定集合 $A$ 中的一个等价关系，用 $\mathcal{E}$ 记由这个关系决定的所有等价类的族。上述引理证明了 $\mathcal{E}$ 中不同元素无交。进一步，由于 $A$ 的每一个元素属于某一等价类中，于是 $\mathcal{E}$ 中元素的并等于整个 $A$。族 $\mathcal{E}$ 就是称之为 $A$ 的分拆的一个特例。

\begin{definition}[分拆]\label{definition:ch01-分拆}
集合 $A$ 的一个分拆（partition）是 $A$ 的无交子集的一个族，其并为 $A$。
\end{definition}

研究集合 $A$ 中的等价关系与研究 $A$ 的分拆实际上是一回事。对于 $A$ 的任意分拆 $\mathcal{D}$，则恰好存在 $A$ 中的一个等价关系，由它可导出 $\mathcal{D}$。

证明这一点并不困难。为了证明 $\mathcal{D}$ 由某一个等价关系导出，我们定义 $A$ 上的关系 $C$ 为：若 $x$ 与 $y$ 属于 $\mathcal{D}$ 的同一元素，则 $xCy$。$C$ 的对称性是显然的。由于 $\mathcal{D}$ 中元素的并等于 $A$，可得到自反性。由于 $\mathcal{D}$ 中不同元素无交得出传递性。容易验证，由 $C$ 决定的等价类的族确实就是族 $\mathcal{D}$。

为了证明这种等价关系的唯一性，设 $C_{1}$ 和 $C_{2}$ 是 $A$ 上的两个等价关系，它们产生出相同的等价类的族 $\mathcal{D}$。对于 $x \in A$，我们证明 $xC_{1}y$ 当且仅当 $xC_{2}y$，从而推得 $C_{1} = C_{2}$。设 $E_{1}$ 是由 $x$ 决定的、相对于关系 $C_{1}$ 的等价类。$E_{2}$ 是由 $x$ 决定的、相对于关系 $C_{2}$ 的等价类。则 $E_{1}$ 是 $\mathcal{D}$ 的一个元素，从而它必定等于 $\mathcal{D}$ 中含有 $x$ 的唯一元素 $D$。类似地，$E_{2}$ 也一定等于 $D$。根据定义，$E_{1}$ 包含所有使得 $yC_{1}x$ 的 $y$，$E_{2}$ 包含所有使得 $yC_{2}x$ 的元素 $y$。因此 $E_{1} = D = E_{2}$，结论得证。

\begin{example}\label{example:ch01-等价类圆}
将平面上的两个点的等价定义为它们与原点的距离相等——它们显然具有自反性、对称性和传递性，则等价类的族 $\mathcal{E}$ 由所有以原点为中心的圆及原点这个单点集组成。
\end{example}

\begin{example}\label{example:ch01-等价类水平线}
将平面上的两个点的等价定义为它们的 $y$ 坐标相同，则等价类的族是平面上所有平行于 $x$ 轴的直线的族。
\end{example}

\begin{example}\label{example:ch01-分拆线族}
设 $\mathcal{L}$ 为平面上与 $y = -x$ 平行的所有直线的族，则 $\mathcal{L}$ 是平面的一个分拆，因为平面的任意点都在一条这样的直线上，并且不同的两条直线无交。分拆 $\mathcal{L}$ 是由平面上的以下等价关系所产生的：如果 $x_{0} + y_{0} = x_{1} + y_{1}$，则点 $(x_{0}, y_{0})$ 与 $(x_{1}, y_{1})$ 等价。
\end{example}

\begin{example}\label{example:ch01-非分拆线族}
设 $\mathcal{L}'$ 为平面上所有直线的族，则 $\mathcal{L}'$ 不是平面的一个分拆，因为 $\mathcal{L}'$ 的不同元素未必无交，两条不重合的直线也可能相交。
\end{example}

\subsection*{序关系}

集合 $A$ 中的一个关系 $C$ 称为序关系（order relation）（或全序（simple order），线序（linear order）），如果满足下列性质：
\begin{enumerate}[(1)]
\item（可比较性）对于 $A$ 中满足 $x \neq y$ 的每个 $x$ 和 $y$，或者 $xCy$，或者 $yCx$.
\item（非自反性）$A$ 中没有 $x$，使得 $xCx$ 成立.
\item（传递性）若 $xCy$ 并且 $yCz$，则 $xCz$.
\end{enumerate}

请注意，性质 (1) 本身并不排除 $A$ 中可能有某元素偶对 $x$，$y$，使 $xCy$ 和 $yCx$ 都成立（因为“或”的含义是“一个或另一个，或者一个及另一个”）。但是与性质 (2) 与 (3) 合在一起，就排除了这种可能；因为如果 $xCy$ 和 $yCx$ 都成立，由传递性得到 $xCx$，这与非自反性矛盾。

\begin{example}\label{example:ch01-通常全序}
考虑由实直线上所有满足 $x < y$ 的实数对 $(x, y)$ 组成的关系。这是一个全序关系，称为实直线上的“通常的全序关系”。下面是一种在实直线上不太为人所熟悉的全序关系：当 $x^{2} < y^{2}$，或者当 $x^{2} = y^{2}$ 并且 $x < y$ 时，有 $xCy$。可以验证这是一种全序关系。
\end{example}

\begin{example}\label{example:ch01-人物关系序}
仍考虑\cref{example:ch01-人物关系}中给出的人与人之间的关系，血缘关系 $B$ 不满足全序关系的任何性质，关系 $S$ 仅仅满足性质 (3)。后代关系比较好，它满足性质 (2) 与 (3)，但仍不满足可比较性。对于在数学中常常出现的满足性质 (2) 和 (3) 的关系，给它一个专门的名字，称之为严格偏序关系。我们以后将研究它（参见第 11 节）。
\end{example}

正如用波纹号 $\sim$ 表示等价关系那样，“小于”符号“$<$”往往用来表示一个全序关系。用这个记号，就可以把一个全序关系的性质写成：
\begin{enumerate}[(1)]
\item 若 $x \neq y$，则或者 $x < y$，或者 $y < x$.
\item 若 $x < y$，则 $x \neq y$.
\item 若 $x < y$ 并且 $y < z$，则 $x < z$.
\end{enumerate}

我们把“或者 $x < y$，或者 $x = y$”论断用记号 $x \leqslant y$ 表示，“$x < y$”论断用 $y > x$ 表示。我们用 $x < y < z$ 表示“$x < y$ 和 $y < z$”。

\begin{definition}[开区间、紧接前元与紧接后元]\label{definition:ch01-开区间}
若 $X$ 是一个集合，$<$ 为 $X$ 上的一个全序关系。对于 $a < b$，我们用记号 $(a, b)$ 表示集合
\begin{align*}
\{x \mid a < x < b\},
\end{align*}
并且称之为 $X$ 中的一个开区间（open interval）。如果这个集合是空集，则称 $a$ 为 $b$ 的紧接前元（immediate predecessor），$b$ 为 $a$ 的紧接后元（immediate successor）。
\end{definition}

\begin{definition}[序型]\label{definition:ch01-序型}
设 $A$ 和 $B$ 是分别有全序关系 $<_{A}$ 和 $<_{B}$ 的两个集合。$A$ 和 $B$ 称为序型（order type）相同的，如果在它们之间有一个一一保序对应，也就是存在一个一一对应 $f: A \longrightarrow B$，使得
\begin{align*}
a_{1} <_{A} a_{2} \Longrightarrow f(a_{1}) <_{B} f(a_{2}).
\end{align*}
\end{definition}

\par
\begin{figure}[htbp]
\centering
\includegraphics[width=0.45\textwidth]{ch01-fig3-2}
\caption{$f(x) = \dfrac{x}{1 - x^{2}}$ 的图形}\label{figure:ch01-保序函数}
\end{figure}

\begin{example}\label{example:ch01-保序对应}
实数区间 $(-1, 1)$ 与实数集 $\mathbb{R}$ 的序型相同，可以验证，由下式
\begin{align*}
f(x) = \frac{x}{1 - x^{2}}
\end{align*}
所给出的函数 $f: (-1, 1) \longrightarrow \mathbb{R}$ 就是一个一一保序对应。函数的图形如\cref{figure:ch01-保序函数}所示。

\end{example}

\begin{example}\label{example:ch01-两段保序}
$\mathbb{R}$ 的两个子集 $A = \{0\} \cup (1, 2)$ 与 $[0, 1) = \{x \mid 0 \leqslant x < 1\}$ 有相同序型。由下式
\begin{align*}
f(0) = 0,
\end{align*}
\begin{align*}
f(x) = x - 1, \quad x \in (1, 2)
\end{align*}
所定义的函数 $f: A \longrightarrow [0, 1)$ 就是所要求的保序对应。
\end{example}

下面是定义全序关系的一种有趣的方式，它在以后的例子中将会用到。

\begin{definition}[字典序]\label{definition:ch01-字典序}
设 $A$ 和 $B$ 是分别有全序关系 $<_{A}$ 和 $<_{B}$ 的两个集合。$A \times B$ 上的全序关系定义为：当 $a_{1} <_{A} a_{2}$ 或者当 $a_{1} = a_{2}$ 并且 $b_{1} <_{B} b_{2}$ 时，
\begin{align*}
a_{1} \times b_{1} < a_{2} \times b_{2}.
\end{align*}
它称为 $A \times B$ 上的字典序关系（dictionary order relation）。
\end{definition}

可以分成几种情形来验证这是一个全序关系，请读者自己完成。

我们选择这个术语的理由是易于理解的，因为定义 $<$ 的法则与在英文字典中字的排序法则是一样的。给定两个英语单词，首先比较它们的首字母，则单词的顺序就可以根据它们的首字母在字母表中出现的次序来决定。如果它们的首字母相同，再比较它们的第二个字母，并且相应地决定其顺序。以此类推。

\begin{example}\label{example:ch01-平面字典序}
考虑平面 $\mathbb{R} \times \mathbb{R}$ 上的字典序。在这个全序关系下，点 $p$ 小于经过 $p$ 的竖直线位于 $p$ 上方的任意点，也小于位于这条竖直线右边的任意点。
\end{example}

\par
\begin{figure}[htbp]
\centering
\includegraphics[width=0.6\textwidth]{ch01-fig3-3}
\caption{两种字典序序型的比较}\label{figure:ch01-字典序比较}
\end{figure}

\begin{example}\label{example:ch01-字典序序型}
考虑具有通常的全序关系的实数集 $[0, 1)$ 及正整数集 $\mathbb{Z}_{+}$，赋予 $\mathbb{Z}_{+} \times [0, 1)$ 字典序，则 $\mathbb{Z}_{+} \times [0, 1)$ 与非负实数集的序型相同，函数
\begin{align*}
f(n \times t) = n + t - 1
\end{align*}
就是所需要的一一保序对应。另一方面，具有字典序的集合 $[0, 1) \times \mathbb{Z}_{+}$ 则有完全不同的序型。例如，这个全序集的每一个元素有一个紧接后元。这些集合如\cref{figure:ch01-字典序比较}所示。

\end{example}

或许你已经知道实数有一个“上确界性质”，对于任意一个全序集也可以定义这个性质。首先，给出下面一些预备定义。

设 $A$ 为具有全序关系 $<$ 的集合。$A_{0}$ 为 $A$ 的一个子集。如果 $b \in A_{0}$ 并且对于任意 $x \in A_{0}$ 有 $x \leqslant b$，则称元素 $b$ 为 $A_{0}$ 的最大元（largest element）。类似地，如果 $a \in A_{0}$ 并且对于任意 $x \in A_{0}$ 有 $a \leqslant x$，则称元素 $a$ 为 $A_{0}$ 的最小元（smallest element）。容易看出，一个集合最多有一个最大元，也最多有一个最小元。

满足以下条件时，我们说 $A$ 的子集 $A_{0}$ 是有上界的（bounded above）：如果存在 $A$ 的一个元素 $b$，对于任意 $x \in A_{0}$，有 $x \leqslant b$，元素 $b$ 称为 $A_{0}$ 的一个上界（upper bound）。如果 $A_{0}$ 的所有上界的集合有一个最小元，则称它为 $A_{0}$ 的上确界（least upper bound，supremum），记作 $\sup A_{0}$。上确界可以属于 $A_{0}$，也可以不属于 $A_{0}$。如果它属于 $A_{0}$，就是 $A_{0}$ 的最大元。

类似地，我们说 $A_{0}$ 是有下界的（bounded below）：如果存在 $A$ 的一个元素 $a$，使得对于任意 $x \in A_{0}$，有 $a \leqslant x$，元素 $a$ 称为 $A_{0}$ 的一个下界（lower bound）。若 $A_{0}$ 的所有下界的集合有一个最大元，则称它为 $A_{0}$ 的下确界（greatest lower bound，infimum），记作 $\inf A_{0}$。下确界可以属于 $A_{0}$，也可以不属于 $A_{0}$。如果它属于 $A_{0}$，就是 $A_{0}$ 的最小元。

现在定义上确界性质。

\begin{definition}[上确界性质]\label{definition:ch01-上确界性质}
如果全序集 $A$ 的每一个有上界的非空子集 $A_{0}$ 必有上确界，则称 $A$ 具有上确界性质（least upper bound property）。同样，如果全序集 $A$ 的每一个有下界的非空子集 $A_{0}$ 必有下确界，则称 $A$ 有下确界性质（greatest lower bound property）。
\end{definition}

作为习题，请读者证明：$A$ 具有上确界性质当且仅当它具有下确界性质。

\begin{example}\label{example:ch01-上确界例}
考虑具有通常的全序关系的实数集 $A = (-1, 1)$。已经假定实数具有上确界性质，那么这个集合具有上确界性质。事实上，任意给定一个 $A$ 的有上界的子集，其上确界（在实数中）必定在 $A$ 中。例如，$A$ 的子集 $\left\{-\dfrac{1}{2n} \mid n \in \mathbb{Z}_{+}\right\}$ 虽然没有最大元，但在 $A$ 中有一个上确界 $0$。

另一方面，集合 $B = (-1, 0) \cup (0, 1)$ 不具有上确界性质。例如，子集 $\left\{-\dfrac{1}{2n} \mid n \in \mathbb{Z}_{+}\right\}$ 以 $(0, 1)$ 中任意元素为其上界，但是它在 $B$ 中没有上确界。
\end{example}

\subsection*{习题}

\subsubsection*{等价关系}

\begin{exercise}\label{exercise:ch01-关系-1}
对于平面的两个点 $(x_{0}, y_{0})$ 和 $(x_{1}, y_{1})$，当 $y_{0} - x_{0}^{2} = y_{1} - x_{1}^{2}$ 时规定它们是等价的。验证这是一个等价关系并给出等价类。
\end{exercise}

\begin{exercise}\label{exercise:ch01-关系-2}
设 $C$ 为集合 $A$ 中的一个关系，$A_{0} \subset A$。$C$ 在 $A_{0}$ 上的限制定义为关系 $C \cap (A_{0} \times A_{0})$。证明等价关系的限制是一个等价关系。
\end{exercise}

\begin{exercise}\label{exercise:ch01-关系-3}
这里给出任何满足对称性和传递性的关系也满足自反性的一个“证明”：“因为 $C$ 有对称性，由 $aCb$ 得出 $bCa$。因为 $C$ 有传递性，由 $aCb$ 及 $bCa$ 得出 $aCa$，证毕。”请找出其中的错误。
\end{exercise}

\begin{exercise}\label{exercise:ch01-关系-4}
设 $f: A \longrightarrow B$ 是一个满射。用
\begin{align*}
f(a_{0}) = f(a_{1})
\end{align*}
定义 $A$ 中的一个关系 $a_{0} \sim a_{1}$。
\begin{enumerate}[(a)]
\item 证明这是一个等价关系。
\item 设 $A^{\ast}$ 为等价类的集合。证明在 $A^{\ast}$ 与 $B$ 之间有一个一一对应。
\end{enumerate}
\end{exercise}

\begin{exercise}\label{exercise:ch01-关系-5}
设 $S$ 和 $S'$ 为平面的两个子集：
\begin{align*}
S = \{(x, y) \mid y = x + 1 \text{ 和 } 0 < x < 2\},
\end{align*}
\begin{align*}
S' = \{(x, y) \mid y - x \text{ 是一个整数}\}.
\end{align*}
\begin{enumerate}[(a)]
\item 证明 $S'$ 是实直线中的一个等价关系并且 $S' \supset S$。给出 $S'$ 的等价类。
\item 证明，对于集合 $A$ 上的等价类的任意一个族，其交也是 $A$ 中的一个等价关系。
\item 给出实直线中的一个等价关系 $T$，使之成为实直线中包含 $S$ 的所有等价关系的交。写出 $T$ 的等价类。
\end{enumerate}
\end{exercise}

\subsubsection*{序关系}

\begin{exercise}\label{exercise:ch01-关系-6}
定义平面中的一个关系为：当 $y_{0} - x_{0}^{2} < y_{1} - x_{1}^{2}$，或当 $y_{0} - x_{0}^{2} = y_{1} - x_{1}^{2}$ 并且 $x_{0} < x_{1}$ 时
\begin{align*}
(x_{0}, y_{0}) < (x_{1}, y_{1}).
\end{align*}
证明这是平面上的一个全序关系，并给出几何解释。
\end{exercise}

\begin{exercise}\label{exercise:ch01-关系-7}
证明：全序关系的限制是一个全序关系。
\end{exercise}

\begin{exercise}\label{exercise:ch01-关系-8}
验证\cref{example:ch01-通常全序}中定义的关系是一个全序关系。
\end{exercise}

\begin{exercise}\label{exercise:ch01-关系-9}
验证字典序是一个全序关系。
\end{exercise}

\begin{exercise}\label{exercise:ch01-关系-10}
(a) 证明\cref{example:ch01-保序对应}中的映射 $f: (-1, 1) \longrightarrow \mathbb{R}$ 是保序的。

(b) 证明由 $g(y) = 2y / [1 + (1 + 4y^{2})^{1/2}]$ 所定义的函数 $g: \mathbb{R} \longrightarrow (-1, 1)$ 既是 $f$ 的左逆又是 $f$ 的右逆。
\end{exercise}

\begin{exercise}\label{exercise:ch01-关系-11}
证明：全序集的一个元素最多有一个紧接后元，也最多有一个紧接前元。证明：全序集的一个子集最多有一个最小元，也最多有一个最大元。
\end{exercise}

\begin{exercise}\label{exercise:ch01-关系-12}
设 $\mathbb{Z}_{+}$ 表示正整数集。考虑 $\mathbb{Z}_{+} \times \mathbb{Z}_{+}$ 上的下列全序关系：
\begin{enumerate}[(i)]
\item 字典序.
\item 当 $x_{0} - y_{0} < x_{1} - y_{1}$，或当 $x_{0} - y_{0} = x_{1} - y_{1}$ 并且 $y_{0} < y_{1}$ 时，$(x_{0}, y_{0}) < (x_{1}, y_{1})$.
\item 当 $x_{0} + y_{0} < x_{1} + y_{1}$，或当 $x_{0} + y_{0} = x_{1} + y_{1}$ 并且 $y_{0} < y_{1}$ 时，$(x_{0}, y_{0}) < (x_{1}, y_{1})$.
\end{enumerate}
在这些全序关系中，什么元素有紧接前元？这个集合有一个最小元吗？证明这三个序型互不相同。
\end{exercise}

\begin{exercise}\label{exercise:ch01-关系-13}
证明下述定理：
\begin{theorem}\label{theorem:ch01-上下确界等价}
如果一个全序集 $A$ 具有上确界性质，则它也具有下确界性质。
\end{theorem}
\end{exercise}

\begin{exercise}\label{exercise:ch01-关系-14}
若 $C$ 为集合 $A$ 中的一个关系，定义 $A$ 上一个新的关系 $D$ 为：当 $(a, b) \in C$ 时，$(b, a) \in D$。
\begin{enumerate}[(a)]
\item 证明：$C$ 有对称性当且仅当 $C = D$.
\item 证明：若 $C$ 是一个全序关系，则 $D$ 也是一个全序关系.
\item 证明习题\cref{exercise:ch01-关系-13}中那个定理的逆定理.
\end{enumerate}
\end{exercise}

\begin{exercise}\label{exercise:ch01-关系-15}
假定实直线具有上确界性质。
\begin{enumerate}[(a)]
\item 证明集合
\begin{align*}
[0, 1] = \{x \mid 0 \leqslant x \leqslant 1\},
\end{align*}
\begin{align*}
[0, 1) = \{x \mid 0 \leqslant x < 1\}
\end{align*}
具有上确界性质.
\item 在字典序下的 $[0, 1] \times [0, 1]$ 具有上确界性质吗？对于 $[0, 1] \times [0, 1)$ 呢？对于 $[0, 1) \times [0, 1]$ 呢？
\end{enumerate}
\end{exercise}

\end{document}